% ===================================================================
%  ਸਾਰਣੀ ਨੰ. 5 — ਘਰ ਤੇ ਉਸ ਦਾ ਬਣਨਾ।
%  Traduction 2026 d'apres texto/io, controlee sur texto/fr, texto/en
%  et texto/hi. Consigne : voir texto/pa/10-sarni-01.tex.
%
%  Le tableau 5 est un DIALOGUE, et les deux volumes composent le nom
%  du parleur en petites capitales : on garde \textsc{}, que le rendu
%  laisse tomber en gardant le texte.
%
%  Les renvois du plan sont des LETTRES — (a) le couloir, (b) le
%  salon... — et non des numeros. On les ecrit comme l'ido les ecrit,
%  y compris la ou l'imprimeur a compose « (1) » pour « (l) ».
%
%  LES DEUX NOMS SONT TRANSCRITS, non traduits, comme dans la colonne
%  shahmoukhi : M. Leblanc est « ਲਬਲਾਂ ਸਾਹਿਬ », M. Roux est
%  « ਰੂ ਸਾਹਿਬ ». Le hindi, lui, traduisait la couleur du premier ; on
%  ne l'a pas suivi, parce qu'une colonne doit se tenir a un seul
%  parti.
% ===================================================================

%%K t05-apar-1 apar
\VUsaut{6.0mm}
\VUtitre{15.0pt}{60}{ਸਾਰਣੀ ਨੰ. 5}
\VUsaut{1.4mm}
\VUtitre{11.4pt}{0}{ਘਰ ਤੇ ਉਸ ਦਾ ਬਣਨਾ।}
\VUsaut{2.2mm}

%%K t05-tit-1 sub
\VUpk{10.2pt}{40}{ਚੰਗੀ ਮੁਲਾਕਾਤ।}
\VUsaut{3.0mm}

%%K t05-01-1 p
\textsc{ਜਾਨ}। --- ਚਾਚਾ ਜੀ, ਮੈਂ ਬਹੁਤ ਜਾਣਨਾ ਚਾਹੁੰਦਾ ਹਾਂ ਕਿ \VUgras{ਘਰ} ਕਿਵੇਂ ਬਣਦਾ
ਹੈ।

%%K t05-01-2 p
\textsc{ਚਾਚਾ} \textsuperscript{(80)}। --- ਇਹ ਸੌਖਾ ਹੈ, ਪੁੱਤਰ; ਮੇਰੇ ਨਾਲ ਆ ਤੇ ਮੈਂ
ਤੈਨੂੰ ਉਹ ਘਰ ਵਿਖਾਵਾਂਗਾ ਜੋ ਬਰਨਾਰਡ ਸਾਹਿਬ \textsuperscript{(79)} ਬਣਵਾ ਰਹੇ ਹਨ ਤੇ
ਜਿਸ ਵਿਚ ਮੈਂ ਰਹਿਣ ਜਾ ਰਿਹਾ ਹਾਂ।

%%K t05-01-3 p
\textsc{ਜਾਨ}। --- ਤੇ ਇਸ ਘਰ ਦਾ \VUgras{ਨਕਸ਼ਾ} \textsuperscript{(76)} ਕਿਸ ਨੇ ਬਣਾਇਆ?

%%K t05-01-4 p
\textsc{ਚਾਚਾ}। --- \VUgras{ਭਵਨ-ਕਾਰ} \textsuperscript{(75)} ਨੇ; ਉਨ੍ਹਾਂ ਨੇ ਬਹੁਤ ਸਾਫ਼
ਨਕਸ਼ਾ ਬਣਾਇਆ, ਉਹ ਮੰਨ ਲਿਆ ਗਿਆ, ਤੇ \VUgras{ਠੇਕੇਦਾਰ} \textsuperscript{(77)} ਨੇ ਕੰਮ
ਸ਼ੁਰੂ ਕਰ ਦਿੱਤਾ।

%%K t05-01-5 p
\textsc{ਜਾਨ}। --- ਤੇ ਠੇਕੇਦਾਰ ਕੌਣ ਹਨ?

%%K t05-01-6 p
\textsc{ਚਾਚਾ}। --- ਲਬਲਾਂ ਸਾਹਿਬ, ਤੇਰੇ ਸਾਥੀ ਜੇਮਜ਼ ਦੇ ਪਿਤਾ। ਉਹ ਭਵਨ-ਕਾਰ ਰੂ ਸਾਹਿਬ
ਨਾਲ ਰਲ ਕੇ ਮਜ਼ਦੂਰਾਂ ਨੂੰ ਹਦਾਇਤਾਂ ਦਿੰਦੇ ਹਨ।

%%K t05-01-7 p
ਲੈ! ਲਬਲਾਂ ਸਾਹਿਬ ਠੀਕ ਵੇਲੇ ਆਪਣਾ \VUgras{ਪੈਮਾਨਾ} \textsuperscript{(78)} ਚੁੱਕੀ ਆ ਰਹੇ
ਹਨ, ਤੇ ਰੂ ਸਾਹਿਬ ਆਪਣਾ ਨਕਸ਼ਾ ਚੁੱਕੀ।

%%K t05-01-8 p
ਨਮਸਕਾਰ, ਜਨਾਬ। ਤੁਸੀਂ ਕਿਵੇਂ ਹੋ?

%%K t05-01-9 p
\textsc{ਰੂ ਸਾਹਿਬ}। --- ਬਹੁਤ ਚੰਗੇ, ਧੰਨਵਾਦ। ਨਮਸਕਾਰ ਜਾਨ; ਮੁੰਡਾ ਕਿੰਨਾ ਵੱਡਾ ਹੋ ਗਿਆ
ਹੈ!

%%K t05-01-10 p
\textsc{ਚਾਚਾ}। --- ਸੱਚੀ ਗੱਲ, ਇਹ ਪੂਰਾ ਨਿੱਕਾ ਬੰਦਾ ਹੈ। ਇਹ ਬਹੁਤ ਗੰਭੀਰ ਹੈ; ਜੋ ਕੁਝ
ਵੇਖਦਾ ਹੈ, ਉਸ ਦਾ ਵੇਰਵਾ ਦੇਣਾ ਪੈਂਦਾ ਹੈ। ਅੱਜ ਇਹ ਜਾਣਨਾ ਚਾਹੁੰਦਾ ਹੈ ਕਿ ਘਰ ਕਿਵੇਂ ਬਣਦਾ
ਹੈ।

%%K t05-tit-2 sub
\VUpk{10.2pt}{40}{ਮਜ਼ਦੂਰ।}
\VUsaut{3.0mm}

%%K t05-01-11 p
\textsc{ਲਬਲਾਂ ਸਾਹਿਬ}। --- ਤੇ ਮੇਰੇ ਨਾਲ ਆ, ਮੇਰੇ ਨਿੱਕੇ ਮਿੱਤਰ, ਤੇ ਮੈਂ ਤੈਨੂੰ ਹੁਣੇ
ਸਮਝਾ ਦਿੰਦਾ ਹਾਂ, ਕਿਉਂਕਿ ਮੈਨੂੰ ਸਿੱਖਣ ਦੇ ਚਾਹਵਾਨ ਬੱਚੇ ਬਹੁਤ ਚੰਗੇ ਲਗਦੇ ਹਨ।

%%K t05-01-12 p
ਉਹ ਮਜ਼ਦੂਰ ਵੇਖਦਾ ਹੈਂ, ਜਿਸ ਦੇ ਹੱਥ ਵਿਚ \VUgras{ਕਹੀ} \textsuperscript{(82)} ਹੈ: ਉਹ
\VUgras{ਪੁਟਾਈ ਕਰਨ ਵਾਲਾ} \textsuperscript{(81)} ਹੈ। ਉਹ ਪਾਣੀ ਦੀ ਨਾਲੀ ਵਿਛਾਉਣ ਲਈ ਇਕ
\VUgras{ਖਾਲ} \textsuperscript{(46)} ਪੁੱਟ ਰਿਹਾ ਹੈ। ਉਸੇ ਨੇ ਉਹ ਜ਼ਮੀਨ ਵੀ ਪੁੱਟੀ ਸੀ
ਜਿੱਥੇ ਘਰ ਖੜਾ ਹੈ, ਤਾਂ ਜੋ ਰਾਜ ਚੌਹਾਂ \VUgras{ਕੰਧਾਂ} ਚੁੱਕ ਸਕੇ। ਉਨ੍ਹਾਂ ਕੰਧਾਂ ਦਾ ਜੋ
ਹਿੱਸਾ ਜ਼ਮੀਨ ਵਿਚ ਹੈ ਉਸ ਨੂੰ \VUgras{ਨੀਂਹ} \textsuperscript{(2)} ਆਖਦੇ ਹਨ। ਨੀਂਹ ਨਾਲ
ਘਿਰੀ ਥਾਂ \VUgras{ਤਹਿਖ਼ਾਨਾ} \textsuperscript{(3)} ਬਣਾਉਂਦੀ ਹੈ। ਇਸ ਤਹਿਖ਼ਾਨੇ ਵਿਚ ਇਕ
ਨਿੱਕੀ ਖਿੜਕੀ ਹੈ ਜਿਸ ਨੂੰ \VUgras{ਰੌਸ਼ਨਦਾਨ} \textsuperscript{(5)} ਆਖਦੇ ਹਨ, ਇਕ
\VUgras{ਬੂਹਾ} \textsuperscript{(4)} ਤੇ ਇਕ ਪੌੜੀ।

%%K t05-01-13 p
\VUgras{ਰਾਜ} \textsuperscript{(108)} ਕੰਧਾਂ ਚੁੱਕਦਾ ਗਿਆ, \VUgras{ਬੂਹਿਆਂ}
\textsuperscript{(8)} ਤੇ \VUgras{ਖਿੜਕੀਆਂ} \textsuperscript{(17)} ਲਈ ਥਾਂ ਛੱਡਦਾ ਹੋਇਆ।

%%K t05-01-14 p
\textsc{ਜਾਨ}। --- ਪਰ ਮੈਨੂੰ ਤਾਂ ਇਹ ਰਾਜ ਵਿਖਾਈ ਹੀ ਨਹੀਂ ਦਿੰਦਾ!

%%K t05-01-15 p
\textsc{ਲਬਲਾਂ ਸਾਹਿਬ}। --- ਉਸ ਦਾ ਘਰ ਦਾ ਕੰਮ ਹੁਣ ਪੂਰਾ ਹੋ ਚੁੱਕਾ ਹੈ। ਉਹ ਉੱਧਰ
\VUgras{ਤਬੇਲੇ} \textsuperscript{(54)} ਕੋਲ ਹੈ, ਆਪਣੇ \VUgras{ਪਾੜ}
\textsuperscript{(110)} ਉੱਤੇ ਚੜ੍ਹਿਆ, \VUgras{ਧੋਬੀ-ਘਰ} \textsuperscript{(56)} ਦੀ ਕੰਧ
ਬਣਾ ਰਿਹਾ ਹੈ।

%%K t05-01-16 p
ਉਸ ਕੋਲ ਇਕ ਕਰਨੀ, ਇਕ ਨਾਂਦ ਤੇ ਇਕ \VUgras{ਸਾਹਲ} \textsuperscript{(109)} ਹੈ। ਪਰ ਇਹ ਨਾਂ
ਤੈਨੂੰ ਯਾਦ ਨਾ ਰਹਿਣਗੇ।

%%K t05-01-17 p
\textsc{ਜਾਨ}। --- ਕਿਉਂ ਨਾ ਰਹਿਣਗੇ, ਮੈਂ ਚਾਚਾ ਜੀ ਕੋਲੋਂ ਇਕ ਨਿੱਕੀ ਕਰਨੀ ਤੇ ਇਕ ਨਾਂਦ
ਮੰਗਾਂਗਾ, ਤੇ ਸਾਹਲ ਤਾਂ ਮੈਂ ਸੁਤਲੀ ਦੇ ਟੁਕੜੇ ਨਾਲ ਆਪ ਹੀ ਬਣਾ ਲਵਾਂਗਾ।

%%K t05-tit-3 sub
\VUsaut{3.0mm}
\VUpk{10.2pt}{40}{ਸਾਮਾਨ।}
\VUsaut{3.0mm}

%%K t05-01-18 p
\textsc{ਚਾਚਾ}। --- ਵਾਹ, ਬਹੁਤ ਚੰਗਾ। ਪਰ ਦੱਸ ਤਾਂ ਜਾਨ, ਘਰ ਬਣਾਉਣ ਲਈ ਕੀ ਚਾਹੀਦਾ ਹੈ?

%%K t05-01-19 p
\textsc{ਜਾਨ}। --- \VUgras{ਤਰਾਸ਼ਿਆ ਹੋਇਆ ਪੱਥਰ} \textsuperscript{(59)} ਤੇ
\VUgras{ਗਾਰਾ} \textsuperscript{(65)} ਚਾਹੀਦਾ ਹੈ।

%%K t05-01-20 p
\textsc{ਚਾਚਾ}। --- ਠੀਕ। ਉਸ ਵੱਡੀ ਹੈਟ ਵਾਲੇ ਬੰਦੇ ਨੂੰ ਵੇਖ, ਆਪਣੇ \VUgras{ਗੱਡੇ}
\textsuperscript{(136)} ਉੱਤੇ। ਉਹ \VUgras{ਗੱਡੇ ਵਾਲਾ} \textsuperscript{(135)} ਹੈ। ਉਹ
ਹਰ ਦਿਨ ਇੱਥੇ \VUgras{ਪੱਥਰ ਦੇ ਢੀਮੇ} \textsuperscript{(60)} ਲਿਆਉਂਦਾ ਹੈ। ਉਹ ਆਪਣਾ ਗੱਡਾ
ਵਿਹੜੇ ਵਿਚ ਖ਼ਾਲੀ ਕਰਦਾ ਹੈ, ਤੇ \VUgras{ਪੱਥਰ ਤਰਾਸ਼} \textsuperscript{(83)} ਢੀਮਿਆਂ ਨੂੰ
ਤਰਾਸ਼ਦੇ ਤੇ ਆਪਣੀ \VUgras{ਪੱਥਰ-ਆਰੀ} \textsuperscript{(85)} ਨਾਲ ਚੀਰਦੇ ਹਨ। ਇਕ
\VUgras{ਮਜ਼ਦੂਰ} \textsuperscript{(146)} ਉਨ੍ਹਾਂ ਨੂੰ ਰਾਜ ਤਕ ਪਹੁੰਚਾਉਂਦਾ ਹੈ, ਜੋ
ਉਨ੍ਹਾਂ ਨੂੰ ਆਪਣੀ ਥਾਂ ਬਿਠਾਉਂਦਾ ਹੈ।

%%K t05-01-21 p
ਇਸ ਦੌਰਾਨ \VUgras{ਗਾਰਾ ਬਣਾਉਣ ਵਾਲਾ} \textsuperscript{(87)} ਗਾਰਾ ਤਿਆਰ ਕਰਦਾ ਹੈ। ਉਸ
ਨੂੰ \VUgras{ਰੇਤ} \textsuperscript{(62)}, \VUgras{ਚੂਨਾ} \textsuperscript{(63)} ਤੇ
\VUgras{ਪਾਣੀ} \textsuperscript{(64)} ਚਾਹੀਦਾ ਹੈ। ਚੂਨਾ ਉਸ ਕੋਲ ਇਕ \VUgras{ਬੋਰੀ}
\textsuperscript{(71)} ਵਿਚ ਹੈ; ਇਕ \VUgras{ਰਾਜ ਦਾ ਸਹਾਇਕ} \textsuperscript{(111)} ਉਸ
ਨੂੰ \VUgras{ਠੇਲੇ} \textsuperscript{(112)} ਵਿਚ ਰੇਤ ਲਿਆਉਂਦਾ ਹੈ, ਤੇ \VUgras{ਸ਼ਾਗਿਰਦ}
\textsuperscript{(113)} ਪੰਪ \textsuperscript{(58)} ਉੱਤੇ ਹੈ। ਉਹ ਇਕ \VUgras{ਬਾਲਟੀ}
\textsuperscript{(114)} ਭਰ ਰਿਹਾ ਹੈ। ਗਾਰਾ ਛੇਤੀ ਹੀ ਤਿਆਰ ਹੋ ਜਾਵੇਗਾ।
\VUgras{ਗਾਰਾ ਚੁੱਕਣ ਵਾਲਾ} \textsuperscript{(107)} ਉਸ ਨੂੰ ਲੈ ਕੇ ਪੌੜੀ ਤੋਂ ਪਾੜ ਉੱਤੇ
ਚੜ੍ਹਾ ਲੈ ਜਾਂਦਾ ਹੈ, ਜਿੱਥੇ ਇਕ ਰਾਜ ਘਰ ਦਾ ਉਹ ਪਾਸਾ ਪੂਰਾ ਕਰ ਰਿਹਾ ਹੈ।

%%K t05-01-22 p
ਤੇ ਹੁਣ ਦੱਸ, ਜੋ ਮਜ਼ਦੂਰ \VUgras{ਛੱਤ} \textsuperscript{(31)} ਬਣਾਉਂਦਾ ਹੈ ਉਸ ਨੂੰ ਕੀ
ਆਖਦੇ ਹਨ?

%%K t05-01-23 p
\textsc{ਜਾਨ}। --- ਉਹ \VUgras{ਛੱਤ ਬਣਾਉਣ ਵਾਲਾ} \textsuperscript{(118)} ਹੈ।

%%K t05-tit-4 sub
\VUsaut{3.0mm}
\VUpk{10.2pt}{40}{ਘਰ ਦੀ ਛੱਤ।}
\VUsaut{3.0mm}

%%K t05-01-24 p
\textsc{ਲਬਲਾਂ ਸਾਹਿਬ}। --- ਹਾਂ, ਪਰ ਪਹਿਲਾਂ \VUgras{ਤਰਖਾਣ} \textsuperscript{(115)}
ਆਉਂਦਾ ਹੈ।

%%K t05-01-25 p
ਤਰਖਾਣ \VUgras{ਛੱਤ ਦਾ ਢਾਂਚਾ} \textsuperscript{(35)} ਬਣਾਉਂਦਾ ਹੈ। ਉਹ
\VUgras{ਕੜੀਆਂ} \textsuperscript{(37)} ਨੂੰ \VUgras{ਕਿੱਲਿਆਂ} \textsuperscript{(117)}
ਨਾਲ ਜੋੜਦਾ ਹੈ। ਉੱਤੇ ਉਹ ਮਜ਼ਦੂਰ, ਆਪਣੀ ਚੌੜੀ ਮਖ਼ਮਲ ਦੀ ਨਿੱਕਰ ਵਿਚ, ਤਰਖਾਣ ਹਨ। ਇਕ ਨਿੱਕੀ
ਕੜੀ ਫੜੀ ਹੋਈ ਹੈ ਤੇ ਦੂਜਾ ਆਪਣੇ \VUgras{ਕੁਹਾੜੇ} \textsuperscript{(116)} ਨਾਲ ਇਕ ਕਿੱਲਾ
ਘੜ ਰਿਹਾ ਹੈ। \VUgras{ਚਿਮਨੀ} \textsuperscript{(41)} ਕੋਲ ਵਾਲਾ ਛੱਤ ਬਣਾਉਣ ਵਾਲਾ ਹੈ। ਉਹ
\VUgras{ਸ਼ਤੀਰਾਂ} (*) ਉੱਤੇ ਪੱਟੀਆਂ ਠੋਕਦਾ ਹੈ, ਤੇ ਪੱਟੀਆਂ ਉੱਤੇ \VUgras{ਸਲੇਟਾਂ}
\textsuperscript{(66)}। ਕਦੇ ਕਦੇ ਉਹ ਖਪਰੈਲਾਂ ਵਿਛਾਉਂਦਾ ਹੈ।

%%K t05-noto-1 noto
\VUnotes{143.2mm}{%
(*) \VUgras{Balk-o} = ਜਰਮਨ \textit{Balken}, ਅੰਗਰੇਜ਼ੀ \textit{joist},
ਫ਼ਰਾਂਸੀਸੀ \textit{solive}, ਇਤਾਲਵੀ \textit{travicello}, ਰੂਸੀ \textit{balka},
ਸਪੇਨੀ ਤੇ ਪੁਰਤਗਾਲੀ \textit{viga}। ਇਕ ਤਕਨੀਕੀ ਧਾਤੂ ਜੋ ਹਾਲੇ ਅਪਣਾਈ ਨਹੀਂ ਗਈ, ਪਰ ਇੱਥੇ
ਫ਼ਰਾਂਸੀਸੀ ਸ਼ਬਦ \textit{solive} ਦਾ ਅਰਥ ਦੇਣ ਲਈ ਸੁਝਾਈ ਜਾਂਦੀ ਹੈ, ਜੋ ਨਾ
\textit{trabo} (ਫ਼ਰਾਂਸੀਸੀ \textit{poutre}) ਹੈ ਨਾ ਨਿੱਕੀ ਕੜੀ। ਉਹ ਚੌਕੋਰ ਘੜੀ ਲੱਕੜ
ਹੈ ਜੋ ਇਕ ਫ਼ਰਸ਼ ਤੇ ਇਕ ਛੱਤ ਸੰਭਾਲਦੀ ਹੈ।%
}

%%K t05-01-26 p
ਤਬੇਲੇ ਦੀ ਛੱਤ \VUgras{ਖਪਰੈਲਾਂ ਦੀ} \textsuperscript{(55)} ਹੈ, ਘਰ ਦੀ
\VUgras{ਸਲੇਟ ਦੀ} \textsuperscript{(32)}, ਤੇ ਬਰਾਂਡੇ ਦੀ \VUgras{ਜਸਤ ਦੀ}
\textsuperscript{(24)}।

%%K t05-01-27 p
\textsc{ਜਾਨ}। --- ਕੀ ਜਸਤ ਦੀਆਂ ਛੱਤਾਂ ਵੀ ਛੱਤ ਬਣਾਉਣ ਵਾਲਾ ਹੀ ਬਣਾਉਂਦਾ ਹੈ?

%%K t05-01-28 p
\textsc{ਲਬਲਾਂ ਸਾਹਿਬ}। --- ਨਹੀਂ, ਉਹ \VUgras{ਜਸਤ-ਸਾਜ਼} \textsuperscript{(121)} ਜਾਂ
\VUgras{ਨਲਸਾਜ਼} ਦਾ ਕੰਮ ਹੈ — ਉਹੋ ਜੋ ਘਰ ਦੀ ਛੱਤ ਵਿਚ ਇਕ \VUgras{ਰੌਸ਼ਨਦਾਨ-ਖਿੜਕੀ}
\textsuperscript{(38)} ਬਿਠਾ ਰਿਹਾ ਹੈ। ਉਹੋ \VUgras{ਛੱਤ ਦੇ ਪਰਨਾਲੇ}
\textsuperscript{(43)} ਤੇ \VUgras{ਉਤਰਨ ਵਾਲੀਆਂ ਨਾਲੀਆਂ} \textsuperscript{(42)} ਵੀ
ਲਾਉਂਦਾ ਹੈ। ਉਹ ਜਸਤ ਤੇ ਟੀਨ ਨੂੰ ਟਾਂਕਾ ਲਾਉਂਦਾ ਹੈ। ਉਸੇ ਨੇ \VUgras{ਬਿਜਲੀ-ਰੋਕ}
\textsuperscript{(40)} ਤੇ \VUgras{ਹਵਾ-ਸੂਚਕ} \textsuperscript{(39)}
\VUgras{ਮੁੰਡੇਰ} \textsuperscript{(36)} ਉੱਤੇ ਗੱਡੇ ਸਨ।

%%K t05-01-29 p
\textsc{ਜਾਨ}। --- ਤੇ ਕੀ ਘਰ ਪੂਰਾ ਹੋ ਗਿਆ?

%%K t05-tit-5 sub
\VUsaut{3.0mm}
\VUpk{10.2pt}{40}{ਔਜ਼ਾਰ।}
\VUsaut{3.0mm}

%%K t05-01-30 p
\textsc{ਲਬਲਾਂ ਸਾਹਿਬ}। --- ਪੂਰਾ ਨਹੀਂ। \VUgras{ਪਲਸਤਰ ਕਰਨ ਵਾਲਾ}
\textsuperscript{(99)} \VUgras{ਇੱਟਾਂ} \textsuperscript{(61)} ਨਾਲ ਉਹ ਕੰਧਾਂ ਬਣਾਉਂਦਾ
ਹੈ ਜੋ ਘਰ ਨੂੰ ਵੱਖੋ-ਵੱਖ ਕਮਰਿਆਂ ਵਿਚ ਵੰਡਦੀਆਂ ਹਨ। ਉਹ \VUgras{ਪਲਸਤਰ}
\textsuperscript{(70)} \VUgras{ਛੱਤਾਂ} \textsuperscript{(26)} ਤੇ ਕੰਧਾਂ ਉੱਤੇ ਚਾੜ੍ਹਦਾ
ਹੈ। ਇਸ ਵੇਲੇ ਉਹ \VUgras{ਬਰਾਂਡੇ} \textsuperscript{(33)} ਦੀ ਛੱਤ ਉੱਤੇ ਕੰਮ ਕਰ ਰਿਹਾ ਹੈ।
ਰਾਜ ਵਾਂਗ ਉਸ ਕੋਲ ਵੀ ਇਕ \VUgras{ਕਰਨੀ} \textsuperscript{(100)} ਤੇ ਇਕ \VUgras{ਨਾਂਦ}
\textsuperscript{(101)} ਹੈ।

%%K t05-01-31 p
ਫਿਰ ਤਰਖਾਣ ਬੂਹੇ, ਖਿੜਕੀਆਂ, \VUgras{ਲੱਕੜ ਦੇ ਫ਼ਰਸ਼} \textsuperscript{(16)} ਤੇ
\VUgras{ਝਿਲਮਿਲਾਂ} \textsuperscript{(19)} ਬਣਾਉਂਦਾ ਹੈ।

%%K t05-01-32 p
ਇਸ ਵੇਲੇ ਉਹ ਵਿਹੜੇ ਵਿਚ ਆਪਣੀ \VUgras{ਕਾਰੀਗਰੀ ਦੀ ਮੇਜ਼} \textsuperscript{(90)} ਉੱਤੇ
ਕੰਮ ਕਰ ਰਿਹਾ ਹੈ। ਉਸ ਕੋਲ ਇਕ \VUgras{ਆਰੀ} \textsuperscript{(94)} ਹੈ, \VUgras{ਲੱਕੜ}
\textsuperscript{(69)} ਚੀਰਨ ਲਈ, ਇਕ \VUgras{ਰੰਦਾ} \textsuperscript{(91)}, ਇਕ
\VUgras{ਸ਼ਿਕੰਜਾ} \textsuperscript{(92)}, ਇਕ \VUgras{ਛੈਣੀ}
\textsuperscript{(94 \textit{bis})}, ਇਕ \VUgras{ਪਰਕਾਰ} \textsuperscript{(95)} ਤੇ
ਲੱਕੜ ਦਾ ਇਕ \VUgras{ਮੁੰਗਲਾ} \textsuperscript{(95 \textit{bis})}।

%%K t05-01-33 p
\textsc{ਜਾਨ}। --- ਓਹੋ, ਤਦੋਂ ਤਾਂ ਤਰਖਾਣ ਹੋਣਾ ਰਾਜ ਹੋਣ ਨਾਲੋਂ ਵੀ ਵੱਧ ਮਜ਼ੇ ਦਾ ਹੈ।
ਚਾਚਾ ਜੀ, ਤੁਸੀਂ ਮੈਨੂੰ ਇਕ ਕਾਰੀਗਰੀ ਦੀ ਮੇਜ਼, ਇਕ ਆਰੀ ਤੇ ਇਕ ਰੰਦਾ ਦਿਵਾ ਦਿਓਗੇ? ਮੈਂ
ਮੇਦੋਰ ਲਈ ਇਕ ਘਰ ਬਣਾਵਾਂਗਾ।

%%K t05-01-34 p
\textsc{ਚਾਚਾ}। --- ਹਾਂ ਪੁੱਤਰ।

%%K t05-01-35 p
\textsc{ਜਾਨ}। --- ਕਿੰਨੀ ਖ਼ੁਸ਼ੀ ਹੋਈ! ਤੇ ਉਹ ਬੰਦਾ ਕੌਣ ਹੈ ਜੋ ਅਗਲੇ ਬੂਹੇ ਕੋਲ ਖੜਾ ਹੈ?

%%K t05-tit-6 sub
\VUsaut{3.0mm}
\VUpk{10.2pt}{40}{ਰੰਗਾਈ।}
\VUsaut{3.0mm}

%%K t05-01-36 p
\textsc{ਲਬਲਾਂ ਸਾਹਿਬ}। --- ਉਹ \VUgras{ਰੰਗਸਾਜ਼} \textsuperscript{(102)} ਹੈ। ਉਹ ਆਪਣੀ
ਪੌੜੀ ਉੱਤੇ ਖੜਾ ਹੈ, \VUgras{ਰੰਗ} \textsuperscript{(103)} ਦੀ ਇਕ ਹਾਂਡੀ ਤੇ ਆਪਣਾ
\VUgras{ਬੁਰਸ਼} \textsuperscript{(104)} ਚੁੱਕੀ। ਉਹ ਅਗਲੇ ਬੂਹੇ ਦੀ ਲੱਕੜ ਉੱਤੇ ਰੰਗ ਚਾੜ੍ਹ
ਰਿਹਾ ਹੈ ਤਾਂ ਜੋ ਉਹ ਵੱਧ ਟਿਕੇ।

%%K t05-01-37 p
ਕਮਰਿਆਂ ਵਿਚ ਕੰਧ ਦਾ ਕਾਗਜ਼ ਵੀ ਉਹੋ ਚਿਪਕਾਉਂਦਾ ਹੈ।

%%K t05-01-38 p
\VUgras{ਸਜਾਵਟ ਕਰਨ ਵਾਲਾ} \textsuperscript{(107)}, ਜੋ \VUgras{ਪੌੜੀ}
\textsuperscript{(106)} ਉੱਤੇ ਚੜ੍ਹਿਆ \VUgras{ਬਾਲਾਖ਼ਾਨੇ} \textsuperscript{(28)} ਉੱਤੇ
ਹੈ, ਇਕ \VUgras{ਪਰਦਾ} \textsuperscript{(29)} ਲਾ ਰਿਹਾ ਹੈ।

%%K t05-01-39 p
ਵੱਡੀ ਖਿੜਕੀ ਕੋਲ ਖੜਾ ਮਜ਼ਦੂਰ \VUgras{ਸ਼ੀਸ਼ਾ-ਸਾਜ਼} \textsuperscript{(96)} ਹੈ। ਉਹ
\VUgras{ਸ਼ੀਸ਼ੇ} \textsuperscript{(18)} \VUgras{ਮਸਾਲੇ} \textsuperscript{(73)} ਨਾਲ
ਜਮਾਉਂਦਾ ਹੈ। ਉਹ ਦੂਜਾ ਮਜ਼ਦੂਰ, ਜੋ ਤਹਿਖ਼ਾਨੇ ਦੇ ਬੂਹੇ ਕੋਲ ਗੋਡਿਆਂ ਭਾਰ ਹੈ,
\VUgras{ਜੰਦਰਾ-ਸਾਜ਼} \textsuperscript{(97)} ਹੈ। ਉਹ ਇਕ \VUgras{ਜੰਦਰਾ}
\textsuperscript{(6)} ਬਿਠਾ ਰਿਹਾ ਹੈ। ਉਸ ਦੀ \VUgras{ਰੇਤੀ} \textsuperscript{(98)} ਉਸ
ਕੋਲ ਪਈ ਹੈ।

%%K t05-01-40 p
\textsc{ਜਾਨ}। --- ਜੋ ਬੰਦਾ ਆਪਣੇ \VUgras{ਹਥੌੜੇ} \textsuperscript{(84)} ਨਾਲ ਲੋਹੇ ਦੇ
ਟੁਕੜੇ ਉੱਤੇ ਸੱਟ ਮਾਰ ਰਿਹਾ ਹੈ, ਕੀ ਉਹ ਵੀ ਜੰਦਰਾ-ਸਾਜ਼ ਹੈ?

%%K t05-01-41 p
\textsc{ਲਬਲਾਂ ਸਾਹਿਬ}। --- ਠੀਕ ਆਖੀਏ ਤਾਂ ਉਹ \VUgras{ਲੁਹਾਰ} \textsuperscript{(125)}
ਹੈ। ਉਹ \VUgras{ਲੋਹੇ ਦਾ ਸਾਮਾਨ} \textsuperscript{(67)} ਘੜ ਰਿਹਾ ਹੈ
\VUgras{ਬਿਜਲੀ ਵਾਲੇ} \textsuperscript{(124)} ਲਈ, ਜੋ \VUgras{ਗੱਡੀ-ਖ਼ਾਨੇ}
\textsuperscript{(51)} ਦੀ ਛੱਤ ਉੱਤੇ ਹੈ। ਉਸ ਕੋਲ ਇਕ \VUgras{ਭੱਠੀ}
\textsuperscript{(126)}, ਇਕ \VUgras{ਆਹਰਣ} \textsuperscript{(127)} ਤੇ ਇਕ
\VUgras{ਸੰਨ੍ਹੀ} \textsuperscript{(128)} ਹੈ।

%%K t05-01-42 p
ਬਿਜਲੀ ਵਾਲਾ ਟੈਲੀਫ਼ੋਨ ਦੀ \VUgras{ਤਾਂਬੇ ਦੀ ਤਾਰ} \textsuperscript{(44)} ਜੋੜ ਰਿਹਾ ਹੈ।

%%K t05-tit-7 sub
\VUpk{10.2pt}{40}{ਬਾਗ਼ਬਾਨੀ।}
\VUsaut{3.0mm}

%%K t05-01-43 p
\textsc{ਚਾਚਾ}। --- \VUgras{ਵਿਹੜੇ} \textsuperscript{(45)} ਦੇ ਸਿਰੇ ਉੱਤੇ,
\VUgras{ਜੰਗਲੇ} \textsuperscript{(47)} ਤੇ \VUgras{ਗੱਡੀ ਦੇ ਫਾਟਕ}
\textsuperscript{(48)} ਕੋਲ, ਤੈਨੂੰ \VUgras{ਮਾਲੀ} \textsuperscript{(129)} ਵਿਖਾਈ ਦਿੰਦਾ
ਹੈ। ਉਹ ਨਿੱਕਾ \VUgras{ਬਾਗ਼} \textsuperscript{(49)} ਤਿਆਰ ਕਰ ਰਿਹਾ ਹੈ। ਉਸ ਨੇ ਆਪਣੀ
\VUgras{ਕਹੀ} \textsuperscript{(131)} ਨਾਲ ਮਿੱਟੀ ਪੁੱਟ ਸੁੱਟੀ ਹੈ, ਉਹ ਬੀ ਬੀਜ ਰਿਹਾ ਹੈ ਤੇ
\VUgras{ਪੰਜੇ} \textsuperscript{(130)} ਨਾਲ ਪੱਧਰਾ ਕਰ ਰਿਹਾ ਹੈ। ਫਿਰ ਆਪਣੇ
\VUgras{ਫੁਹਾਰੇ} \textsuperscript{(132)} ਨਾਲ ਉਹ \VUgras{ਕਿਆਰੀਆਂ}
\textsuperscript{(150)} ਨੂੰ ਸਿੰਜੇਗਾ ਤੇ ਬੂਟੇ ਉੱਗਣਗੇ।

%%K t05-01-44 p
\VUgras{ਸਾਈਸ} \textsuperscript{(133)}, ਜੋ ਹੁਣੇ ਘੋੜੇ ਦੀ ਸੰਭਾਲ ਕਰ ਕੇ ਉਸ ਨੂੰ ਖੁਆ
ਚੁੱਕਾ ਹੈ, \VUgras{ਬੱਘੀ} \textsuperscript{(52)} ਸਾਫ਼ ਕਰ ਰਿਹਾ ਹੈ — ਸਪੰਜ, ਇਕ
\VUgras{ਹੱਥ-ਪੰਪ} \textsuperscript{(134)} ਤੇ ਪਾਣੀ ਦੀ ਬਾਲਟੀ ਨਾਲ। ਫਿਰ ਉਹ ਉਸ ਨੂੰ
ਗੱਡੀ-ਖ਼ਾਨੇ ਵਿਚ ਰੱਖ ਦੇਵੇਗਾ ਤੇ ਭਾਰੇ \VUgras{ਅਰਗਲ} \textsuperscript{(53)} ਨਾਲ ਬੂਹਾ
ਬੰਦ ਕਰ ਦੇਵੇਗਾ।

%%K t05-01-45 p
ਲੈ, ਹੁਣ ਤਾਂ ਤਸੱਲੀ ਹੋਈ ਹੋਣੀ; ਸਮਝਾਉਣਾ ਬਹੁਤ ਹੋ ਚੁੱਕਾ।

%%K t05-01-46 p
\textsc{ਜਾਨ}। --- ਹਾਂ ਚਾਚਾ ਜੀ, ਪਰ ਮੈਂ ਘਰ ਦੇ ਅੰਦਰ ਵੀ ਵੇਖਣਾ ਚਾਹਾਂਗਾ।

%%K t05-01-47 p
\textsc{ਚਾਚਾ}। --- ਉਹ ਕੱਲ੍ਹ ਲਈ। ਰੂ ਸਾਹਿਬ ਤੈਨੂੰ ਨਕਸ਼ਾ ਸਮਝਾਉਣਗੇ ਤੇ ਹਰ ਕਮਰੇ ਦਾ ਨਾਂ
ਦੱਸਣਗੇ।

%%K t05-tit-8 sub
\VUsaut{3.0mm}
\VUpk{10.2pt}{40}{ਘਰ ਦੀ ਅੰਦਰਲੀ ਬਣਤਰ।}
\VUsaut{2.4mm}

%%K t05-apar-2 apar
{\centering\textit{(ਨਕਸ਼ਾ ਵੇਖੋ।)}\par}
\VUsaut{2.4mm}

%%K t05-01-48 p
\textsc{ਭਵਨ-ਕਾਰ}। --- ਆ ਜਾਨ, ਮੈਂ ਤੈਨੂੰ ਘਰ ਦੇ ਵੱਖੋ-ਵੱਖ ਹਿੱਸੇ ਵਿਖਾਉਂਦਾ ਹਾਂ।

%%K t05-01-49 p
ਚੱਲ \VUgras{ਹੇਠਲੀ ਮੰਜ਼ਿਲ} \textsuperscript{(7)} ਵਿਚ ਚੱਲੀਏ। ਇਹ ਰਿਹਾ
\VUgras{ਘੰਟੀ} ਦਾ ਬਟਨ \textsuperscript{(11)}। ਅਸੀਂ ਤਿੰਨ \VUgras{ਪੌੜੀਆਂ}
\textsuperscript{(14)} ਚੜ੍ਹਦੇ ਹਾਂ, ਜੋ \VUgras{ਡਿਉੜ੍ਹੀ} \textsuperscript{(9)} ਦੀਆਂ
ਹਨ, ਫਿਰ \VUgras{ਦਹਿਲੀਜ਼} \textsuperscript{(10)} ਲੰਘਦੇ ਹਾਂ, ਜੋ \VUgras{ਅਗਲੇ ਬੂਹੇ}
\textsuperscript{(8)} ਦੀ ਹੈ, ਤੇ ਇਹ ਰਹੇ ਅਸੀਂ \VUgras{ਪੌੜੀਆਂ ਦੇ ਥੜੇ}
\textsuperscript{(13)} ਦੇ ਹੇਠਾਂ।

%%K t05-01-50 p
\textsc{ਜਾਨ}। --- ਇਹ ਲਾਂਘਾ ਹੈ।

%%K t05-01-51 p
\textsc{ਭਵਨ-ਕਾਰ}। --- ਨਹੀਂ, ਇਹ \VUgras{ਡਿਉੜ੍ਹੀ-ਕਮਰਾ} \textsuperscript{(12)} ਹੈ।
\VUgras{ਲਾਂਘਾ} \textsuperscript{(a)} ਸਿਰੇ ਉੱਤੇ ਹੈ। ਸੱਜੇ ਪਾਸੇ ਇਹ ਰਹੇ
\VUgras{ਬੈਠਕ} \textsuperscript{(b)} ਤੇ \VUgras{ਖਾਣੇ ਦਾ ਕਮਰਾ} \textsuperscript{(c)};
ਖਾਣੇ ਦੇ ਕਮਰੇ ਦੇ ਸਾਹਮਣੇ \VUgras{ਰਸੋਈ} \textsuperscript{(d)} ਤੇ
\VUgras{ਭੰਡਾਰ-ਕਮਰਾ} \textsuperscript{(e)}; ਫਿਰ ਅਗਲੇ ਪਾਸੇ
\VUgras{ਪੜ੍ਹਨ ਦਾ ਕਮਰਾ} \textsuperscript{(f)}। \VUgras{ਬਰਾਂਡਾ}
\textsuperscript{(23)}, ਆਪਣੇ \VUgras{ਸ਼ੀਸ਼ੇ ਦੇ ਬੂਹੇ} \textsuperscript{(25)} ਸਮੇਤ,
ਬੈਠਕ ਕੋਲ ਹੈ।

%%K t05-01-52 p
ਹੁਣ ਅਸੀਂ ਪੌੜੀਆਂ ਚੜ੍ਹਾਂਗੇ। \VUgras{ਜੰਗਲਾ} \textsuperscript{(15)} ਫੜੀ ਰੱਖ ਤੇ
ਪੌੜੀਆਂ ਗਿਣ। ਚੌਦਾਂ ਹਨ। ਇਹ ਰਹੀ \VUgras{ਚੌਕੀ} \textsuperscript{(m)}: ਅਸੀਂ ਪਹਿਲੀ
\VUgras{ਮੰਜ਼ਿਲ} \textsuperscript{(27)} ਉੱਤੇ ਹਾਂ।

%%K t05-tit-9 sub
\VUsaut{3.0mm}
\VUpk{10.2pt}{40}{ਪਹਿਲੀ ਮੰਜ਼ਿਲ।}
\VUsaut{3.0mm}

%%K t05-01-53 p
ਪੌੜੀਆਂ ਦੇ ਸਾਹਮਣੇ \VUgras{ਪਖ਼ਾਨਾ} \textsuperscript{(20)} ਤੇ
\VUgras{ਕੱਪੜਿਆਂ ਦਾ ਕਮਰਾ} \textsuperscript{(j)} ਹਨ। ਸੱਜੇ ਪਾਸੇ ਇਕ ਸੋਹਣਾ
\VUgras{ਸੌਣ ਦਾ ਕਮਰਾ} \textsuperscript{(g)}, ਜਿਸ ਦੀਆਂ ਦੋ ਖਿੜਕੀਆਂ ਘਰ ਦੇ
\VUgras{ਅਗਲੇ ਰੁਖ਼} \textsuperscript{(1)} ਵੱਲ ਖੁੱਲ੍ਹਦੀਆਂ ਹਨ। ਖੱਬੇ ਪਾਸੇ
\VUgras{ਨ੍ਹਾਉਣ ਦਾ ਕਮਰਾ} \textsuperscript{(1)} ਤੇ \VUgras{ਪ੍ਰਾਹੁਣਿਆਂ ਦਾ ਕਮਰਾ}
\textsuperscript{(i)} ਹਨ, ਜਿਸ ਵਿਚ ਇਕ ਵੱਡਾ \VUgras{ਆਲਾ} \textsuperscript{(h)} ਹੈ।
ਸੌਣ ਦੇ ਕਮਰੇ ਕੋਲ \VUgras{ਬੱਚਿਆਂ ਦਾ ਕਮਰਾ} \textsuperscript{(k)} ਹੈ। ਉਹ ਇਕ ਚੌੜੇ
\VUgras{ਛੱਤ ਦੇ ਬਾਰਜੇ} \textsuperscript{(21)} ਵਿਚ ਖੁੱਲ੍ਹਦਾ ਹੈ। ਇਸ ਬਾਰਜੇ ਦੇ ਹੇਠਾਂ ਇਕ
ਹੋਰ ਪਖ਼ਾਨਾ \textsuperscript{(20)} ਹੈ, ਤੇ ਕੰਧ ਨਾਲ ਲੱਗੀ ਇਹ ਰਹੀ \VUgras{ਘੰਟੀ}
\textsuperscript{(22)}, ਜੋ ਖਾਣੇ ਲਈ ਵਜਾਈ ਜਾਂਦੀ ਹੈ।

%%K t05-01-54 p
\textsc{ਜਾਨ}। --- ਚੱਲੋ \VUgras{ਦੂਜੀ ਮੰਜ਼ਿਲ} ਉੱਤੇ ਚੱਲੀਏ।

%%K t05-01-55 p
\textsc{ਭਵਨ-ਕਾਰ}। --- ਦੂਜੀ ਮੰਜ਼ਿਲ ਹੈ ਹੀ ਨਹੀਂ। ਛੱਤ ਦੇ ਹੇਠਾਂ ਸਿਰਫ਼
\VUgras{ਚੁਬਾਰੇ} \textsuperscript{(33)} ਤੇ ਭੰਡਾਰ \textsuperscript{(34)} ਹਨ। ਸਾਡਾ
ਚੱਕਰ ਪੂਰਾ ਹੋਇਆ, ਹੁਣ ਹੇਠਾਂ ਉਤਰ ਸਕਦੇ ਹਾਂ।

%%K t05-01-56 p
\textsc{ਜਾਨ}। --- ਧੰਨਵਾਦ ਜਨਾਬ, ਬਹੁਤ ਦਿਲਚਸਪ ਰਿਹਾ। ਮੈਂ ਆਪਣੇ ਪਿਤਾ ਨੂੰ ਸਭ ਕੁਝ
ਦੱਸਾਂਗਾ ਜੋ ਮੈਂ ਵੇਖਿਆ।
\VUsaut{4.0mm}
\VUfilet{22mm}
